% =====================================================================
\section{Датчики: температура, влажность, расстояние}
% =====================================================================
\lessonintro
  {подключать цифровые датчики температуры, влажности и расстояния и уметь читать описание (datasheet) незнакомого датчика}
  {урок 3 (подтяжка), урок 6 (делитель напряжения), урок 1 (опрос по \code{millis()} без пауз)}
  {плата измеряет температуру, влажность и расстояние и печатает результаты в Serial Monitor}
  {DS18B20 и резистор \SI{4.7}{k\ohm}; DHT22; HC-SR04 и резисторы \SI{1}{k\ohm} и \SI{2}{k\ohm}}

\subsection{Сначала: датчик --- это собеседник со своими правилами}

Цифровой датчик сам измеряет температуру или расстояние и присылает плате готовое число. Но у каждого
датчика свои правила разговора: какое напряжение ему нужно, по какому проводу и в каком порядке
передавать сигналы, как часто можно спрашивать. Если нарушить правила, датчик либо промолчит, либо
ответит ерунду, а иногда и выйдет из строя. Все правила записаны в datasheet --- паспорте датчика. Этот
урок учит находить в нём главное на примере трёх популярных датчиков.

\subsection{Как читать datasheet за пять минут}

Не нужно читать документ целиком. Для первого подключения достаточно найти пять вещей.

\begin{enumerate}
  \item \textbf{Питание.} Допустимый диапазон напряжения (например, \SIrange{3}{5.5}{V}) и потребляемый ток.
  \item \textbf{Логические уровни.} Каким напряжением датчик «говорит». Если выходной сигнал выше \SI{3.3}{V},
        вход ESP32-S3 нужно защищать.
  \item \textbf{Интерфейс.} Один провод, I\textsuperscript{2}C, SPI или импульсы; есть ли подтяжка.
  \item \textbf{Время и период опроса.} Сколько длится замер и как часто можно его повторять.
  \item \textbf{Предельные значения} (Absolute Maximum Ratings). Выход за них --- и датчик необратимо портится.
\end{enumerate}

\begin{center}\small
\begin{tabularx}{\textwidth}{@{}l X X X@{}}
\toprule
 & DS18B20 & DHT22 & HC-SR04 \\
\midrule
Что измеряет & температура & температура и влажность & расстояние \\
Питание & \SIrange{3}{5.5}{V} & \SIrange{3.3}{6}{V} & \SI{5}{V} \\
Сигнал & по уровню питания & по уровню питания & ECHO даёт \SI{5}{V} \\
Интерфейс & 1-Wire & свой однопроводной & импульсы TRIG и ECHO \\
Период опроса & замер до \SI{750}{ms} & не чаще раза в \SI{2}{s} & не чаще раза в \SI{60}{ms} \\
Диапазон & от $-55$ до $+125$\,°C & \SIrange{0}{100}{\percent} влажности & \SIrange{2}{400}{cm} \\
\bottomrule
\end{tabularx}
\end{center}

Значения в таблице ориентировочные: у разных производителей и версий модулей они немного отличаются.
Всегда сверяйтесь с документом на свой экземпляр.

\subsection{DS18B20: температура по шине 1-Wire}

DS18B20 --- маленький датчик температуры в корпусе, похожем на транзистор, или в герметичном щупе. Он
использует шину \textbf{1-Wire}: один провод данных плюс питание и земля. Выход датчика
--- с открытым стоком, как у I\textsuperscript{2}C (урок 7), поэтому линию данных к питанию притягивает
\emph{подтягивающий резистор} \SI{4.7}{k\ohm}. Без него шина не работает.

У каждого DS18B20 есть уникальный 64-битный адрес, записанный на заводе. Поэтому несколько датчиков
подключают к одному и тому же выводу платы и различают по адресам.

\begin{figure}[H]
\centering
\begin{circuitikz}[scale=0.95,transform shape]
  \node[blk,fill=blkcpu,minimum width=2.8cm,minimum height=3.8cm] (m) at (0,0) {};
  \node[blk,fill=blkper,minimum width=2.8cm,minimum height=3.8cm] (d) at (8,0) {};
  \node[font=\sffamily\small\bfseries] at (0,2.3) {ESP32-S3};
  \node[font=\sffamily\small\bfseries,align=center] at (8,2.45) {DS18B20\\[-2pt]\scriptsize вид снизу};
  \foreach \y/\n in {1.5/3V3,0/{\pin{7}},-1.5/GND}
    \node[font=\sffamily\scriptsize,anchor=east] at (1.4,\y) {\n};
  \foreach \y/\n in {1.5/VDD,0/DQ,-1.5/GND}
    \node[font=\sffamily\scriptsize,anchor=west] at (6.6,\y) {\n};
  \draw[thick,esp]    (1.4,1.5)  -- (6.6,1.5);
  \draw[thick,accent] (1.4,0)    -- (6.6,0);
  \draw[thick,black!70] (1.4,-1.5) -- (6.6,-1.5);
  \draw[thick,esp] (4,1.5) to[R,l=\SI{4.7}{k\ohm},-*] (4,0);
  \fill[esp] (4,1.5) circle (1.6pt);
  \node[font=\sffamily\scriptsize,text=accent] at (4,-0.45) {линия данных};
\end{circuitikz}
\caption{Подключение DS18B20: линия данных одним проводом, подтяжка \SI{4.7}{k\ohm} к 3V3.
Все датчики шины включаются параллельно}
\end{figure}

\codecap{DS18B20: чтение температуры}
\codeinput{l15_ds18b20}

Установите через Library Manager библиотеки \textbf{OneWire} и \textbf{DallasTemperature}. Значение
\code{DEVICE\_DISCONNECTED\_C} (в библиотеке это $-127$) означает, что датчик не ответил. Не принимайте его
за температуру: поставьте проверку, как в примере.

\begin{tip}
Щуп на проводах обычно имеет красный провод (VDD), чёрный (GND) и жёлтый или белый (DQ). Цвета не
стандартизованы, поэтому проверьте маркировку у продавца. Если щуп сильно греется в первые секунды, то
питание и земля перепутаны: сразу отключите его.
\end{tip}

\subsection{DHT22: температура и влажность}

DHT22 (AM2302) измеряет температуру и относительную влажность. Он тоже общается по одному проводу, но
\emph{это не 1-Wire}: у него собственный протокол из коротких импульсов, где ноль и единица
различаются длительностью, измеряемой десятками микросекунд. Поэтому библиотека на время чтения
занимает процессор и не терпит пауз.

Датчик внутри медленный: новые данные у него появляются примерно раз в две секунды, а каждый обмен
слегка его нагревает. Поэтому опрашивать DHT22 чаще раза в \SI{2}{s} нельзя --- вы получите повтор
старого значения или ошибку. В примере для этого стоит проверка по \code{millis()}, а не \code{delay()},
как в уроке 1.

\begin{figure}[H]
\centering
\begin{circuitikz}[scale=0.95,transform shape]
  \node[blk,fill=blkcpu,minimum width=2.8cm,minimum height=4.6cm] (m) at (0,0) {};
  \node[blk,fill=blkper,minimum width=2.8cm,minimum height=4.6cm] (d) at (8,0) {};
  \node[font=\sffamily\small\bfseries] at (0,2.7) {ESP32-S3};
  \node[font=\sffamily\small\bfseries,align=center] at (8,2.85) {DHT22\\[-2pt]\scriptsize 4 ножки, вид спереди};
  \foreach \y/\n in {1.8/3V3,0.6/{\pin{15}},-1.8/GND}
    \node[font=\sffamily\scriptsize,anchor=east] at (1.4,\y) {\n};
  \foreach \y/\n in {1.8/{1~VCC},0.6/{2~DATA},-0.6/{3~NC},-1.8/{4~GND}}
    \node[font=\sffamily\scriptsize,anchor=west] at (6.6,\y) {\n};
  \draw[thick,wred]   (1.4,1.8)  -- (6.6,1.8);
  \draw[thick,wgreen] (1.4,0.6)  -- (6.6,0.6);
  \draw[thick,wblack] (1.4,-1.8) -- (6.6,-1.8);
  \draw[thick,wred] (4,1.8) to[R,l=\SI{10}{k\ohm},-*] (4,0.6);
  \fill[wred] (4,1.8) circle (1.6pt);
  \node[font=\sffamily\scriptsize,text=wgreen!70!black] at (4,0.2) {линия данных};
\end{circuitikz}
\caption{Подключение DHT22 (датчик без платы): данные одним проводом на \pin{15}, подтяжка
\SI{10}{k\ohm} к 3V3. В модуле с тремя контактами (VCC, DATA, GND) резистор уже установлен}
\end{figure}

\codecap{DHT22: температура и влажность раз в 2,5 секунды}
\codeinput{l15_dht22}

В Library Manager нужны \textbf{DHT sensor library} и \textbf{Adafruit Unified Sensor}. Если у вас
модуль на плате с тремя контактами, подтяжка на нём уже есть; для «голого» датчика с четырьмя ногами
поставьте резистор \SI{10}{k\ohm} между DATA и 3V3.

\subsection{HC-SR04: расстояние по времени эха}

Ультразвуковой дальномер HC-SR04 посылает звуковой щелчок частотой около \SI{40}{kHz} и слушает эхо.
Чем дальше предмет, тем дольше звук идёт туда и обратно. Плата не получает «число сантиметров» ---
она измеряет \emph{длительность} импульса на выводе ECHO.

\begin{figure}[H]
\centering
\begin{tikzpicture}[x=0.8cm,y=0.75cm,font=\sffamily\scriptsize]
  % TRIG
  \node[left,text=accent] at (-0.3,4.5) {TRIG};
  \draw[very thick,accent] (0,4) -- (1.5,4) -- (1.5,5) -- (2.5,5) -- (2.5,4) -- (12,4);
  \draw[decorate,decoration={brace,amplitude=3pt}] (1.5,5.1) -- (2.5,5.1)
    node[midway,above=3pt,font=\sffamily\scriptsize]{$\geq\SI{10}{\micro s}$};
  % ультразвуковая пачка
  \node[left,text=okgreen,align=right] at (-0.3,2.5) {40 кГц\\[-2pt] (звук)};
  \draw[very thick,okgreen] (0,2) -- (2.8,2)
    \foreach \k in {0,...,7}{ -- ({2.8+\k*0.28},3) -- ({2.94+\k*0.28},3) -- ({2.94+\k*0.28},2) -- ({3.08+\k*0.28},2) }
    -- (12,2);
  \node[text=okgreen,anchor=south west] at (2.7,3.1) {8 импульсов};
  % ECHO
  \node[left,text=esp] at (-0.3,0.5) {ECHO};
  \draw[very thick,esp] (0,0) -- (5.2,0) -- (5.2,1) -- (9.4,1) -- (9.4,0) -- (12,0);
  \draw[decorate,decoration={brace,mirror,amplitude=3pt}] (5.2,-0.2) -- (9.4,-0.2)
    node[midway,below=4pt,font=\sffamily\scriptsize]{$t$ --- время «туда и обратно»};
  % звук вернулся
  \draw[dashed,gray] (9.4,-0.1) -- (9.4,3.3);
  \node[gray,anchor=south west,align=left] at (9.5,2.2) {эхо\\[-2pt] вернулось};
  \draw[dashed,gray] (5.2,-0.1) -- (5.2,1.6);
  \draw[->,>=Stealth,gray] (0,-1.6) -- (12.3,-1.6) node[right]{время};
\end{tikzpicture}
\caption{Измерение расстояния HC-SR04. Импульс на TRIG запускает измерение, датчик излучает пачку звука,
и на ECHO появляется HIGH, длительность которого равна времени пути звука туда и обратно}
\end{figure}

Скорость звука в воздухе около \SI{343}{m/s}, то есть $\SI{0.0343}{cm/\micro s}$. Звук проходит
расстояние дважды, поэтому
\[
  L = \frac{t \cdot 0{,}0343}{2} \;[\text{см}],
\]
где $t$ --- длительность ECHO в микросекундах. Например, при $t = \SI{1166}{\micro s}$ получаем
$L \approx \SI{20}{cm}$.

Измерить длительность импульса можно функцией \code{pulseIn(pin, HIGH, timeout)}. Третий аргумент ---
тайм-аут в микросекундах: если эхо не вернулось (предмет слишком далеко или слишком мягкий), функция
вернёт 0, а не зависнет. Без тайм-аута по умолчанию ждёт целую секунду.

\begin{caution}
Вывод ECHO даёт сигнал \SI{5}{V}, а входы ESP32-S3 рассчитаны на \SI{3.3}{V}. Подключать ECHO к плате
\textbf{напрямую нельзя}: вход может выйти из строя. Поставьте делитель напряжения из резисторов \SI{1}{k\ohm} и
\SI{2}{k\ohm}, как на рисунке ниже.
\end{caution}

\begin{figure}[H]
\centering
\begin{circuitikz}[scale=0.95,transform shape]
  \node[blk,fill=blkcpu,minimum width=2.8cm,minimum height=6.4cm] (m) at (0,0) {};
  \node[blk,fill=blkper,minimum width=2.8cm,minimum height=6.4cm] (d) at (8.6,0) {};
  \node[font=\sffamily\small\bfseries] at (0,3.8) {ESP32-S3};
  \node[font=\sffamily\small\bfseries,align=center] at (8.6,3.95) {HC-SR04};
  \foreach \y/\n in {2.7/5V,0.9/{\pin{16}},-0.9/{\pin{21}},-2.7/GND}
    \node[font=\sffamily\scriptsize,anchor=east] at (1.4,\y) {\n};
  \foreach \y/\n in {2.7/VCC,0.9/TRIG,-0.9/ECHO,-2.7/GND}
    \node[font=\sffamily\scriptsize,anchor=west] at (7.2,\y) {\n};
  \draw[thick,wred]    (1.4,2.7)  -- (7.2,2.7);
  \draw[thick,wgreen]  (1.4,0.9)  -- (7.2,0.9);
  \draw[thick,wblack]  (1.4,-2.7) -- (7.2,-2.7);
  % ECHO через делитель
  \draw[thick,worange] (7.2,-0.9) to[R,l=\SI{1}{k\ohm}] (4.6,-0.9) -- (1.4,-0.9);
  \draw[thick,worange] (3.0,-0.9) to[R,l_=\SI{2}{k\ohm},-*] (3.0,-2.7);
  \fill[worange] (3.0,-0.9) circle (1.6pt);
  \fill[wblack] (3.0,-2.7) circle (1.6pt);
  \node[font=\sffamily\scriptsize,text=worange!70!black,anchor=south] at (6.6,-0.85) {5\,В};
  \node[font=\sffamily\scriptsize,text=worange!70!black,anchor=north west] at (1.5,-1.0) {3,3\,В};
  \node[font=\sffamily\small,anchor=west,text=accent] at (3.6,-1.85)
    {$U = 5\cdot\dfrac{2}{1+2}\approx\SI{3.3}{V}$};
\end{circuitikz}
\caption{Подключение HC-SR04. Питание \SI{5}{V}, TRIG идёт на \pin{16} напрямую, а ECHO попадает на \pin{21}
только через делитель: нижний резистор вдвое больше верхнего, поэтому на вход платы приходит две трети
от \SI{5}{V}}
\end{figure}

Питание самого HC-SR04 --- \SI{5}{V} (вывод 5V платы, который при питании от USB даёт напряжение
шины). TRIG подключается напрямую: плата выдаёт \SI{3.3}{V}, а датчик такой уровень воспринимает как
единицу.

\codecap{HC-SR04: расстояние с тайм-аутом}
\codeinput{l15_hcsr04}

\subsection{Самостоятельный опыт: парковочный радар}

Соберите сигнализатор, как в автомобиле: чем ближе предмет, тем чаще мигает светодиод на \pin{4}. Сделайте
отдельный скетч и разберите задачу на шаги:
\begin{enumerate}
  \item Измеряйте расстояние раз в \SI{100}{ms} функцией из примера.
  \item Переведите расстояние в период мигания: \SI{100}{cm} и дальше --- без мигания, \SI{5}{cm} --- период \SI{50}{ms}
        (подсказка: \code{map()} и \code{constrain()}).
  \item Мигайте через \code{millis()}, а не через \code{delay()}, иначе измерения будут «залипать».
\end{enumerate}
Проверьте работу на плоском твёрдом предмете (книга, ладонь) и на мягком (свитер). У мягких предметов эхо слабее,
поэтому дальность падает.

\subsection{То же на ESPUI: дальномер на экране телефона}

Парковочный радар из самостоятельного опыта удобно смотреть на экране: круговая шкала даёт
мгновенное «далеко/близко», график показывает, как менялось расстояние. Схема подключения
HC-SR04 прежняя (ECHO через делитель!), добавляется только телефон. Программу соберём в три шага.

\subsubsection{Шаг 1. Только датчик}

Сначала убедитесь, что сам дальномер работает, и не трогайте сеть. Функция \code{readDistanceCm()}
оформляет измерение из листинга выше, а результат пока уходит в Serial Monitor.

\codecap{Шаг 1. Расстояние в Serial Monitor}
\codeinput{l15_espui_radar_s1}

\begin{itemize}[leftmargin=*]
  \item \code{readDistanceCm()} возвращает $-1$, если эхо не вернулось, --- ошибка не маскируется под ноль;
  \item замер запускается по \code{millis()} раз в \SI{200}{ms}, то есть не чаще пяти раз в секунду;
  \item при $-1$ печатаем «нет эха», иначе расстояние в сантиметрах.
\end{itemize}

\subsubsection{Шаг 2. Шкала на экране телефона}

Теперь плата сама создаёт Wi-Fi-сеть, а расстояние показывает круговая шкала \code{Gauge}.

\codecap{Шаг 2. Шкала Gauge в панели ESPUI}
\codeinput{l15_espui_radar_s2}

\begin{itemize}[leftmargin=*]
  \item \code{WiFi.softAP(имя, пароль)} --- плата создаёт свою сеть (как в уроке~5);
  \item \code{addControl(ControlType::Gauge, \ldots)} создаёт шкалу, а два элемента \code{Min} и \code{Max}
        задают её края: от 0 до \SI{200}{cm};
  \item \code{updateControlValue(gaugeId, \ldots)} двигает стрелку; \code{constrain()} не даёт ей выйти за шкалу;
  \item неудачные замеры пока просто пропускаются. Это мы исправим на следующем шаге.
\end{itemize}

\subsubsection{Шаг 3. Подпись, график и честная ошибка}

Добавим подпись состояния и график истории. Главное изменение: при ошибке датчика выводится сообщение,
а стрелка и график не получают ложный «нуль».

\codecap{Шаг 3. ESPUI: расстояние HC-SR04 как Gauge и Graph}
\codeinput{l15_espui_radar}

\begin{itemize}[leftmargin=*]
  \item \code{ESPUI.label(\ldots)} создаёт подпись; \code{updateLabel()} меняет её текст;
  \item \code{ESPUI.graph(\ldots)} создаёт график, а \code{addGraphPoint()} добавляет в него очередную точку;
  \item ветка \code{if (d < 0)} показывает «нет эха» вместо расстояния.
\end{itemize}

Подключитесь к сети \code{S3-Radar} (пароль \code{lesson-39-demo}). Поднесите ладонь к датчику:
стрелка уйдёт к нулю, а в графике появится «ступенька». Закройте датчик мягкой тканью или
отведите далеко: вместо нулевого значения на странице появится сообщение «нет эха»,
потому что ошибка датчика не должна выглядеть как «расстояние 0 см».

\begin{trouble}
\begin{itemize}[leftmargin=*]
  \item \textbf{DS18B20 показывает $-127$ °C} --- нет подтяжки \SI{4.7}{k\ohm}, оборван DQ или перепутаны питание и земля.
  \item \textbf{DHT22 выдаёт NaN или повторяет значения} --- опрос чаще раза в 2 с, нет подтяжки или выбран не тот тип датчика в конструкторе.
  \item \textbf{HC-SR04 всегда «нет эха» или показывает случайные числа} --- питание не \SI{5}{V}, перепутаны TRIG и ECHO или нет общей земли.
\end{itemize}
\end{trouble}

\begin{summary}
\begin{itemize}[leftmargin=*]
  \item Перед подключением датчика найдите в datasheet питание, логические уровни, интерфейс и период опроса.
  \item Для DS18B20 и DHT22 нужна подтяжка линии данных, а опрашивать их надо не чаще, чем разрешает документ.
  \item Выход \SI{5}{V} (ECHO у HC-SR04) подключают к ESP32-S3 только через делитель, а время ожидания эха ограничивают тайм-аутом.
\end{itemize}
\end{summary}

\begin{selfcheck}
\begin{enumerate}
  \item Какие пять вещей надо найти в datasheet нового датчика?
  \item Почему нельзя опрашивать DHT22 чаще одного раза в две секунды?
  \item Зачем нужен делитель напряжения для ECHO и как посчитать расстояние по времени эха?
\end{enumerate}
\end{selfcheck}

\begin{tasks}
\begin{enumerate}
\item \lv{1} Прочитайте datasheet своего датчика и выпишите рабочее напряжение, диапазон и точность.
\item \lv{1} Выводите температуру DS18B20 раз в секунду в Serial Monitor.
\item \lv{2} Включайте светодиод на \pin{4}, когда температура выше заданной, с гистерезисом в 1\textdegree{}C.
\item \lv{2} Измерьте расстояние HC-SR04 пять раз на одной дистанции и найдите разброс показаний.
\item \lv{3} Сделайте «электронную линейку»: расстояние выводится на OLED, кнопка запоминает ноль.
\item \lv{4} Подключите к одной линии два DS18B20 и выведите температуру каждого с его адресом.
\end{enumerate}
\end{tasks}
